\documentclass[11pt,a4paper]{article}
\usepackage[utf8]{inputenc}
\usepackage[T1]{fontenc}
\usepackage[romanian]{babel}
\usepackage[margin=2cm]{geometry}
\usepackage{array}
\usepackage{tabularx}
\usepackage{enumitem}
\usepackage{graphicx}
\usepackage{setspace}
\usepackage{xcolor}
\usepackage{hyperref}
\usepackage{amssymb}
\setlength{\parindent}{0pt}
\newcommand{\fillline}[1][6cm]{\underline{\hspace{#1}}}
\newcommand{\fillbox}[1][0.4cm]{\framebox[#1]{\rule{0pt}{#1}}}
\newcommand{\checkboxempty}{$\square$}

\begin{document}

\begin{minipage}{0.3\textwidth}
\textbf{PRIMĂRIA}\\
\textbf{CLUJ-NAPOCA}
\end{minipage}%
\begin{minipage}{0.4\textwidth}
\centering
\textbf{*431001*}
\end{minipage}%
\begin{minipage}{0.3\textwidth}
\raggedleft
Nr. \fillline[3cm] / \fillline[1.5cm]
\end{minipage}

\vspace{0.3cm}
\hrule
\vspace{0.5cm}

\begin{center}
\textbf{Către,}\\
\textbf{PRIMARUL MUNICIPIULUI CLUJ-NAPOCA}\\[0.3cm]
\textbf{C E R E R E}\\
\textbf{PENTRU EMITEREA CERTIFICATULUI DE URBANISM}
\end{center}

\vspace{0.4cm}

Subsemnatul/a$^{1}$ \underline{ {{nume_complet}} },

CNP \underline{ {{cnp}} }, cu$^{2}$ \checkboxempty\ domiciliul / \checkboxempty\ sediul în judeţul \fillline[5cm]

\fillline[3cm], municipiul/oraşul/comuna \fillline[5cm] satul \fillline[4cm],

sectorul \fillline[1.5cm], cod poştal \fillline[2cm], str. \fillline[6cm] nr. \fillline[1.5cm]

\fillline[2cm], bl. \fillline[1cm], sc. \fillline[1cm], et. \fillline[1cm], ap. \fillline[1cm], telefon/fax \fillline[4cm],

e-mail \fillline[5cm] \fillline[5cm] în calitate de /reprezentant

al \fillline[4cm] CUI \fillline[4cm] în conformitate cu prevederile Legii nr. 50/1991,

privind autorizarea executării lucrărilor de construcţii, republicată, cu modificările şi completările ulterioare, solicit emiterea certificatului de urbanism în scopul \underline{ {{scopul_lucrarii}} }

\hrulefill

\hrulefill

\vspace{0.4cm}

\begin{tabularx}{\textwidth}{|X|X|}
\hline
\checkboxempty 1. Elaborarea documentaţiei pentru autorizarea executării lucrărilor de construcţii, în conformitate cu prevederile art.3 alin.(1) din Lege, privind:\newline
\checkboxempty 1.1. Lucrări de construire\newline
\checkboxempty 1.2. Lucrări de desfiinţare\newline
\checkboxempty a) lucrări de construire, reconstruire, consolidare, modificare, extindere, reabilitare, schimbare de destinaţie sau de reparare a construcţiilor de orice fel, precum şi a instalaţiilor aferente acestora, cu excepţia celor prevăzute la art.12 din Legea nr.50/1991;\newline
\checkboxempty b) lucrări de construire, reconstruire, extindere, reparare, consolidare, protejare, restaurare, conservare, precum şi orice alte lucrări, indiferent de valoarea lor, care urmează să fie efectuate la construcţii reprezentând monumente istorice, stabilite potrivit legii;\newline
\checkboxempty c) lucrări de construire, reconstruire, modificare, extindere, reparare, modernizare şi reabilitare privind căile de comunicaţie de orice fel, drumurile forestiere, lucrările de artă, reţelele şi dotările tehnico-edilitare, lucrările hidrotehnice, amenajările de albii, lucrările de îmbunătăţiri funciare, lucrările de instalaţii de infrastructură, lucrările pentru noi capacităţi de producere, transport, distribuţie a energiei electrice şi/sau termice, precum şi de reabilitare şi retehnologizare a celor existente;\newline
\checkboxempty d) împrejmuiri şi mobilier urban, amenajări de spaţii verzi, parcuri, pieţe şi alte lucrări de amenajare a spaţiilor publice;
&
\checkboxempty e) lucrări de foraje si excavări necesare pentru efectuarea studiilor geotehnice si a prospectiunilor geologice, proiectarea si deschiderea exploatarilor de cariere si balastiere, a sondelor de gasze si petrol, precum si a altor exploatari de suprafata sau subterane.\newline
\checkboxempty f) lucrări, amenajări şi construcţii cu caracter provizoriu, necesare în vederea organizării executării lucrărilor, în condiţiile prevăzute la art.7 alin.(11) din Legea nr.50/1991;\newline
\checkboxempty g) organizarea de tabere de corturi, căsuţe sau rulote;\newline
\checkboxempty h) lucrări de construcţii cu caracter provizoriu, chioşcuri, tonete, cabine, spaţii de expunere situate pe căile şi spaţiile publice, corpuri şi panouri de afişaj, firme şi reclame, precum şi anexele gospodăreşti ale exploataţiilor agricole situate în condiţiile prevăzute la art.7 alin.(11) din Legea nr.50/1991;\newline
\checkboxempty i) cimitire noi şi extinderi.\newline
\textbf{\checkboxempty 2. Operaţiuni notariale privind circulaţia imobiliară:}\newline
\checkboxempty vânzări, \checkboxempty cumpărări, \checkboxempty concesionări,\newline
\checkboxempty cesionări, \checkboxempty dezmembrări, \checkboxempty parcelări,\newline
\checkboxempty comasări, \checkboxempty partaje, \checkboxempty succesiuni etc.\newline
\textbf{\checkboxempty 3. Adjudecarea prin licitaţie a proiectării lucrărilor publice (denumire):}\newline
\hrulefill\newline
\hrulefill\newline
\checkboxempty 4. Cereri în justiţie\newline
\checkboxempty 5. Alte scopuri prevăzute de lege (definire)\newline
\hrulefill\\
\hline
\end{tabularx}

\vfill
\hrule
\vspace{0.2cm}
\small
www.primariaclujnapoca.ro \\
telefon: 0264/596.030 \\
Cluj-Napoca, str. Moţilor, nr. 3 \hfill \textbf{*431001*}

\vspace{0.2cm}
\footnotesize
Timp estimativ de completare: 15 minute\\
Datele personale care vă sunt solicitate prin prezenta cerere vor fi prelucrate numai în vederea procesării și soluționării solicitării dumneavoastră. Primăria municipiului Cluj-Napoca garantează securitatea procesării datelor și arhivarea acestora în conformitate cu prevederile legale în vigoare. Responsabilul Primăriei municipiului Cluj-Napoca cu protecția datelor poate fi contactat pe adresa de email dpo@primariaclujnapoca.ro.\\
În conformitate cu Regulamentul nr. 679/2016 aveți dreptul de a solicita Primăriei municipiului Cluj-Napoca, în ceea ce priveşte datele cu caracter personal referitoare la persoana vizată, accesul la acestea, rectificarea sau ştergerea acestora sau restricţionarea prelucrării sau a dreptului de a vă opune prelucrării, precum şi a dreptului la portabilitatea datelor.

\newpage

\begin{minipage}{0.3\textwidth}
\textbf{PRIMĂRIA}\\
\textbf{CLUJ-NAPOCA}
\end{minipage}%
\begin{minipage}{0.4\textwidth}
\centering
\textbf{*431001*}
\end{minipage}%
\begin{minipage}{0.3\textwidth}
\raggedleft
Nr. \fillline[3cm] / \fillline[1.5cm]
\end{minipage}

\vspace{0.3cm}
\hrule
\vspace{0.5cm}

pentru imobilul \checkboxempty\ teren şi/sau \checkboxempty\ construcţii situat în \underline{ {{adresa_imobil}} }

municipiul/oraşul/comuna \fillline[4cm] satul \fillline[4cm], sectorul \fillline[3cm],

cod poştal \fillline[3cm] str. \fillline[6cm] \fillline[3cm],

nr. \fillline[2cm], bl. \fillline[1cm] sc. \fillline[1cm] et. \fillline[1cm], ap. \fillline[1cm], sau identificat prin$^{3}$ \underline{ {{nr_cadastral}} }

\hrulefill

\vspace{0.4cm}

\textbf{În sprijinul identificării imobilului anexez:} planul cadastral/topografic actualizat la zi, scara 1: \fillline[3cm], precum şi extrasul de Carte Funciară pentru informare, eliberate de OCPI$^{4}$;

\hrulefill

\vspace{0.3cm}

\hspace{0.5cm}\textbf{Suprafaţa terenului şi/sau construcţiei pentru care solicit certificatul de urbanism este de \underline{ {{suprafata_constructie}} } m$^{2}$.}

\vspace{0.5cm}

\checkboxempty\ Prin prezenta solicit comunicarea răspunsului pe următoarea adresă de e-mail:\\
\hrulefill

\vspace{0.2cm}

\checkboxempty\ Mă oblig să comunic instituției orice modificare intervine în legătură cu această adresă de e-mail

\checkboxempty\ Îmi exprim consimțământul ca Primăria Municipiului Cluj-Napoca să comunice orice informații, date personale, clarificări și completări pe adresa de e-mail indicată mai sus

\checkboxempty\ Am luat la cunoștință faptul că în cazul nefuncționării serverului de e-mail comunicat sau în cazul adresei greșite de e-mail, Municipiul Cluj-Napoca nu poate fi tras la răspundere pentru acest lucru

\vspace{1cm}

Data: \fillline[3cm] \hfill L.S. \hfill Semnătura$^{5}$ \underline{ {{nume_complet}} }

\vfill
\hrule
\vspace{0.2cm}
\small
www.primariaclujnapoca.ro \\
telefon: 0264/596.030 \\
Cluj-Napoca, str. Moţilor, nr. 3 \hfill \textbf{*431001*}

\newpage

\begin{minipage}{0.3\textwidth}
\textbf{PRIMĂRIA}\\
\textbf{CLUJ-NAPOCA}
\end{minipage}%
\begin{minipage}{0.4\textwidth}
\centering
\textbf{*431001*}
\end{minipage}%
\begin{minipage}{0.3\textwidth}
\raggedleft
Nr. \fillline[3cm] / \fillline[1.5cm]
\end{minipage}

\vspace{0.3cm}
\hrule
\vspace{0.5cm}

\begin{center}
\textbf{P R E C I Z Ă R I}\\
\textbf{privind completarea formularului}\\
\textbf{,,Cerere pentru emiterea certificatului de urbanism''}
\end{center}

\vspace{0.4cm}

\textbf{1) Numele şi prenumele solicitantului:}

\checkboxempty\ persoană fizică; sau

\checkboxempty\ reprezentant al firmei (persoană juridică), cu precizarea denumirii acesteia, precum şi a calităţii solicitantului în cadrul firmei.

\vspace{0.3cm}

\textbf{2) Domiciliul/sediul firmei:}

Pentru persoană fizică: \checkboxempty\ se completează cu date privind domiciliul acesteia.\\
Pentru persoană juridică: \checkboxempty\ se completează cu date privind sediul social al firmei.

\vspace{0.3cm}

\textbf{3) Alte elemente de identificare:}

În situaţia în care amplasamentul imobilului nu este evidenţiat în planurile cadastrale sau topografice ale localităţii/teritoriului administrativ (la scările 1:500, 1:2000 sau 1:10.000, după caz) aflate în gestiunea oficiului de cadastru şi publicitate imobiliară teritorial, pentru identificarea imobilului solicitantul va putea prezenta, după caz, informaţii privind:

\checkboxempty\ localitatea, numărul cadastral şi numărul de carte funciară, în cazul în care legea nu dispune altfel; sau\\
\checkboxempty\ elemente de reper, general cunoscute; sau\\
\checkboxempty\ numărul de ordine şi suprafaţa de teren înscrise în Registrul agricol; sau\\
\checkboxempty\ plan de situaţie extras din cadrul unor studii şi/sau planuri urbanistice elaborate anterior în zonă.

\vspace{0.3cm}

\textbf{4) Planuri cadastrale/topografice, cu evidenţierea imobilelor în cauză, astfel:}

\checkboxempty\ pentru imobilele neînscrise în evidenţele de cadastru şi publicitate imobiliară: Plan de încadrare în zonă, la una din scările 1:10.000, 1:5.000, 1:2.000, 1:1.000, 1:500, după caz, eliberat, la cerere, de către oficiul de cadastru şi publicitate imobiliară;

\checkboxempty\ pentru imobilele înscrise în evidenţele de cadastru şi publicitate imobiliară: Extras din planul cadastral de pe ortofotoplan şi extrasul de carte funciară pentru informare actualizat la zi, eliberate, la cerere, de către oficiul de cadastru şi publicitate imobiliară.

\newpage

\begin{minipage}{0.3\textwidth}
\textbf{PRIMĂRIA}\\
\textbf{CLUJ-NAPOCA}
\end{minipage}%
\begin{minipage}{0.4\textwidth}
\centering
\textbf{*431001*}
\end{minipage}%
\begin{minipage}{0.3\textwidth}
\raggedleft
Nr. \fillline[3cm] / \fillline[1.5cm]
\end{minipage}

\vspace{0.3cm}
\hrule
\vspace{0.5cm}

\textbf{5) La rubrica ,,Semnătura'':}

\checkboxempty\ se va înscrie şi în clar numele solicitantului: persoană fizică sau reprezentant al persoanei juridice.

\vspace{0.4cm}
\hrule
\vspace{0.4cm}

În conformitate cu Legea nr. 50/1991 actualizată art. 6, alin (4) şi Ordinul M.D.R.L nr. 839/2009 pentru aprobarea Normelor Metodologice de aplicare a Legii nr. 50/1991 privind autorizarea executării lucrărilor de construcţii, cap. II, art. 19, alin (1) pentru emiterea certificatului de urbanism, solicitantul trebuie să depună la emitent o documentaţie \textbf{în dublu exemplar}, cuprinzând:

\vspace{0.3cm}

\textbf{a) cerere tip completată la toate rubricile cu:}\\
- elemente de identificare solicitant (nume, adresă, CNP)\\
- elemente de identificare imobil (construcţie sau teren) pentru care se solicită certificat (localitate, stradă, nr. cadastral, nr. Carte funciară)

\vspace{0.2cm}

\textbf{b) planuri cadastrale/topografice}, împreună cu inventarul de coordonate (întocmit şi semnat de topograf autorizat), cu evidenţierea imobiliară în cauză astfel:

\textbf{1. pentru imobilele neînscrise în evidenţele de cadastru şi publicitate imobiliară}\\
plan de încadrare în zona la una din scările 1:10000, 1:5000, după caz, eliberat la cerere de către oficiul de cadastru şi Publicitate imobiliară -- \textbf{în original}

\textbf{2. pentru imobilele înscrise în evidenţele de cadastru şi publicitate imobiliară:} extras din planul cadastral de pe ortofotoplan şi extras de carte funciară pentru informare actualizat la zi, eliberate, la cerere, de către biroul de cadastru şi publicitate imobiliară -- \textbf{în original}

\textbf{c)} extras de carte funciară

\textbf{d)} taxă pentru emiterea certificatului de urbanism, calculată potrivit hotărârii de consiliu local.

\vspace{0.5cm}
\footnotesize
$^{1}$Primăria municipiului Cluj-Napoca asigură gratuit fotocopierea documentelor.\\
Pentru documentele emise de către alte instituții publice, organe de specialitate ale administrației publice centrale și locale, precum și persoane juridice de drept privat, care potrivit legii au obținut statut de utilitate publică sau sunt autorizate să presteze un serviciu public, în regim de putere publică, pe care solicitantul nu le depune, acesta are posibilitatea de a-și da consimțământul expres ca Primăria municipiului Cluj-Napoca să obțină în numele său copii sau extrase ale acestora.

\vfill
\hrule
\vspace{0.2cm}
\small
www.primariaclujnapoca.ro \\
telefon: 0264/596.030 \\
Cluj-Napoca, str. Moţilor, nr. 3 \hfill \textbf{*431001*}

\end{document}
